\documentclass{ctexart}
\usepackage[a4paper,left=3cm,right=3cm,top=3cm,bottom=3cm]{geometry}
\usepackage{amssymb,amsfonts,amsmath,amsthm}
\usepackage{anyfontsize}
\usepackage{graphicx}
\usepackage{tikz-cd}
\usepackage{bm}
\usepackage{bbm}
\usepackage{mathrsfs}
\usepackage[colorlinks = true,      % 启用彩色链接而非带框链接
            linkcolor = blue,       % 内部链接（如目录、交叉引用）颜色
            urlcolor = blue,        % 外部链接颜色
            citecolor = blue]{hyperref} % 参考文献链接颜色

\newtheorem{definition}{Definition}[section]
\newtheorem{remark}{Remark}[section]
\newtheorem{theorem}{Theorem}[section]
\newtheorem{proposition}{Proposition}[section]
\newtheorem{setting}{Setting}[section]
\newtheorem{lemma}{Lemma}[section]
\newtheorem{reason}{Reason}[section]
\newtheorem{example}{Example}[section]
\newtheorem{corollary}{Corollary}[section]
\newtheorem{exercise}{Exercise}[section]

\renewcommand{\proofname}{\textit{Proof}}
\renewcommand{\contentsname}{Contents}
\renewcommand{\refname}{References}
\renewcommand{\hom}{\operatorname{Hom}}

\newcommand{\leftquotes}{``}
\newcommand{\rightquotes}{''}

\newcommand{\mapsfrom}{\mathrel{\reflectbox{\ensuremath{\mapsto}}}}
\newcommand{\longmapsfrom}{\mathrel{\reflectbox{\ensuremath{\longmapsto}}}}

\newcommand{\rank}{\operatorname{rank}}
\newcommand{\im}{\operatorname{im}}
\newcommand{\identity}{\operatorname{id}}
\newcommand{\trdeg}{\operatorname{trdeg}}
\newcommand{\codim}{\operatorname{codim}}
\newcommand{\depth}{\operatorname{depth}}
\newcommand{\height}{\operatorname{height}}
\newcommand{\projdim}{\dim_{\operatorname{proj}}}
\newcommand{\characteristic}{\operatorname{Char}}

\newcommand{\derivation}{\operatorname{Der}}
\newcommand{\ext}{\operatorname{Ext}}
\newcommand{\mor}{\operatorname{Mor}}
\newcommand{\aut}{\operatorname{Aut}}
\newcommand{\homend}{\operatorname{End}}

\newcommand{\sheafhom}{\mathcal{H}\kern -.5pt om}
\newcommand{\sheaftor}{\mathcal{T}\kern -.5pt or}
\newcommand{\sheafext}{\mathcal{E}\kern -.5pt xt}
\newcommand{\sheafend}{\mathcal{E}\kern -.5pt nd}

\newcommand{\coker}{\operatorname{coker}}
\newcommand{\coim}{\operatorname{coim}}
\newcommand{\cone}{\operatorname{cone}}
\newcommand{\exceptionallocus}{\operatorname{Ex}}
\newcommand{\totalcomplex}{\operatorname{Tot}}

\newcommand{\spec}{\operatorname{Spec}}
\newcommand{\proj}{\operatorname{Proj}}
\newcommand{\spf}{\operatorname{Spf}}
\newcommand{\fraction}{\operatorname{Frac}}
\newcommand{\singular}{\operatorname{Sing}}
\newcommand{\support}{\operatorname{Supp}}
\newcommand{\framebundle}{\operatorname{Fr}}

\newcommand{\reduction}{\operatorname{red}}
\newcommand{\normalization}{\operatorname{nor}}
\newcommand{\regular}{\operatorname{reg}}
\newcommand{\smooth}{\operatorname{sm}}
\newcommand{\analytification}{\operatorname{an}}

\newcommand{\inverselimit}{\underset{\longleftarrow}{\lim}} 
\newcommand{\colimit}{\underset{\longrightarrow}{\lim}}
\newcommand{\derivedotimes}{\overset{\llcorner}{\otimes}}

\newcommand{\grassmannianscheme}[2]{\operatorname{Gr}(#1, #2)}
\newcommand{\grassmanmianfunctor}[2]{\mathcal{G}\kern -.5pt r(#1, #2)}
\newcommand{\hilbertscheme}{\operatorname{Hilb}}
\newcommand{\hilbertfunctor}{\mathcal{H}\kern -.5pt ilb}
\newcommand{\picardscheme}{\operatorname{Pic}}
\newcommand{\picardfunctor}{\mathcal{P}\kern -.5pt ic}
\newcommand{\isomorphismscheme}{\operatorname{Isom}}
\newcommand{\isomorphismfunctor}{\mathcal{I}\kern -.5pt som}

\newcommand{\abeliangroups}{\mathbf{Ab}}
\newcommand{\rings}{\mathbf{Ring}}
\newcommand{\commutativerings}{\mathbf{CRing}}
\newcommand{\sets}{\mathbf{Set}}

\newcommand{\module}{\mathbf{Mod}}
\newcommand{\vectorspace}{\mathbf{Vect}}
\newcommand{\algebra}{\mathbf{Alg}}
\newcommand{\scheme}{\mathbf{Sch}}
\newcommand{\vectorbundle}{\underline{\mathbf{Vect}}}

\newcommand{\matrices}[2]{\operatorname{Mat}_{#1}(#2)}
\newcommand{\generallineargroup}[2]{\operatorname{GL}_{#1}(#2)}
\newcommand{\speciallineargroup}[2]{\operatorname{SL}_{#1}(#2)}
\newcommand{\projectivespeciallineargroup}[2]{\operatorname{PSL}_{#1}(#2)}

\newcommand{\order}{\operatorname{ord}}
\newcommand{\divisor}{\operatorname{div}}
\newcommand{\annihilator}{\operatorname{ann}}

\title{Vector Bundles on Curves}
\author{Taught by Jiabin Du}
\date{}

\begin{document}
\maketitle
\newpage

\tableofcontents
\newpage

\section*{Introduction}
\label{section:Introduction}

The course aims at introducing vector bundles on curve and Koszul cohomology.
You may benefit from the following references
\begin{itemize}
    \item \emph{Koszul Cohomology and Algebraic Geometry}, Aprodu-Nagel
    \item \emph{Lectures on Vector Bundles}, J.Le Potier
    \item \emph{The Geometry of Moduli Space of Sheaves}, Huybrechts-Lehn.
\end{itemize}

\section{Vector Bundles}
\label{section:Vector_Bundles}

\subsection{Sheaf theory}
\label{subsection:Vector_Bundles_Sheaf_theory}

Recall the definition of sheaves.
\begin{definition}
    Let $X$ be a topological space and $\mathcal{C}$ be a category.
    Assume $op(X)$ is the category of open subsets of X with morphisms are inclusions.
    A presheaf on $X$ valued on $\mathcal{C}$ (or say of $\mathcal{C}$) is a contravariant functor $\mathcal{F}: op(X) \longrightarrow \mathcal{C}$.
\end{definition}
\begin{remark}
    Though we can consider this general definition of presheaf, we would only consider presheaves at least valued on $Ab$,
    where $Ab$ is the category of abelian groups.
    In addition, we need these presheaves map $\varnothing$ to $\{0\}$.
\end{remark}
\begin{definition}
    Let $X$ be a topological space, $\mathcal{F}$ presheaf of abelian groups on $X$.
    For $V \subseteq V$, we call the correspondence homomorphism $\mathcal{F}(U) \longrightarrow \mathcal{F}(V)$ a restriction map, denoted by $\cdot \rho_{UV}$.
    Elements in $\mathcal{F}(U)$ is called a section of $\mathcal{F}$ on $U$.
    In particular, if $U = X$, we call elements of $\mathcal{F}(X)$ global sections, otherwise local sections.
\end{definition}
\begin{remark}
    For convenience, for $s\in \mathcal{F}(U)$, we would write $s\big{|}_V$ to denote image of $s$ instead of $\rho_{UV}(s)$.
    Also, sometimes $\mathcal{F}(U)$ can be rewritten as $\Gamma(U, \mathcal{F})$.
\end{remark}
\begin{definition}
    Let $X$ be a topological space, $\mathcal{F}$ presheaf of abelian groups on $X$.
    $\mathcal{F}$ is a sheaf on $X$ if it satisfies that \\
    (1)For each open subset $U$ with open covering $\{U_{i}\}_{i\in I}$ and $s\in \mathcal{F}(U)$, 
    if $s\big{|}_{U_{i}} = 0$ for all $i$, then $s = 0$. \\
    (2)For each open subset $U$ with open covering $\{U_{i}\}_{i\in I}$,
    if there exists $s_{i}\in \mathcal{F}(U_{i})$ for all $i$, satisfying $s_{i}\big{|}_{U_{i} \cap U_{j}} = s_{j}\big{|}_{U_{i} \cap U_{j}}$ for all $i,j$,
    then there exists $s\in \mathcal{F}(U)$ such that $s\big{|}_{U_{i}} = s_{i}$ for all $i$.
\end{definition}
\begin{remark}
    The two conditions in fact give a exact sequence for each open subset $U$ with open covering $\{U_{i}\}_{i\in I}$.
    \begin{equation*}
        0 \longrightarrow \mathcal{F}(U) \longrightarrow \prod_{i\in I} \mathcal{F}(U_{i}) \longrightarrow \prod_{i,j\in I} \mathcal{F}(U_{i} \cap U_{j})
    \end{equation*}
\end{remark}
\begin{example}
    (1)Let $X$ be a topological space.
    Then $U \longmapsto \{\text{continuous functions $U \longrightarrow \mathbb{R}$}\}$ give a sheaf on $X$. \\
    (2)Let $X = \mathbb{R}^{n}$.
    Then $U \longmapsto \{\text{$C^{\infty}$-functions $U \longrightarrow \mathbb{R}$}\}$ gives a sheaf on $X$. \\
    (3)Let $X$ be a topological space and $A$ be an abelian group.
    View $A$ as a topological space with discrete topology.
    Then $U \longmapsto \{\text{locally constant functions $U \longrightarrow A$}\}$ gives a sheaf on $X$, which is called the constant sheaf of $A$ on $X$.
\end{example}
\begin{definition}
    Let $X$ be a topological space, $\mathcal{F}$ (pre)sheaf of abelian groups on $X$.
    For $x\in X$, define the stalk of $\mathcal{F}$ at $x$ to be $\mathcal{F}_{x} = \underset{\underset{U\ni x}{\longrightarrow}}{\lim} \mathcal{F}(U)$.
    More concretely, an element of $\mathcal{F}_{x}$ is represented by a pair $(U, s) = s_{x}$ where $U$ is an open neighbourhood of x and $s\in \mathcal{F}(U)$, called a germ.
    Two germs $(U, s)$ and $(V, t)$ represent the same element if there exists open neighbourhood $W$ of $x$ such that $W \subseteq U \cap V$ and $s\big{|}_{W} = t\big{|}_{W}$.
\end{definition}
\begin{proposition}
    Let $X$ be a topological space, $\mathcal{F}$ sheaf of abelian groups on $X$.
    Suppose $U$ is an open subset of $X$ and $s,t\in \mathcal{F}(U)$.
    Then $s = t$ if and only if $s_{x} = t_{x}$ for all $x\in U$.
    \label{Proposition 1.1}
\end{proposition}
\begin{proof}
    $\leftquotes \Rightarrow \rightquotes$: Obviously.
    \par
    $\leftquotes \Leftarrow \rightquotes$: Since for all $x\in U$, we have $s_{x} = t_{x}$,
    there exists open neighbourhood $W_{x}$ of $x$ such that $W_{x} \subseteq U$ and $s\big{|}_{W_{x}} = t\big{|}_{W_{x}}$.
    Note that $\{W_{x}\}_{x\in U}$ form an open covering of $U$.
    By property of sheaf, we get $s = t$.
\end{proof}
\begin{definition}
    Let $X$ be a topological space, $\mathcal{F}, \mathcal{G}$ (pre)sheaves of abelian groups on $X$.
    A morphism $\varphi$ from $\mathcal{F}$ to $\mathcal{G}$ is a family of group homomorphisms $\varphi(U): \mathcal{F}(U) \longrightarrow \mathcal{G}(U)$ such that for $V \subseteq U$,
    there is a commutative diagram
    \begin{equation*}
        \begin{tikzcd}
            \mathcal{F}(U) \arrow[r, "\varphi(U)"] \arrow[d, "\rho_{UV}"] &
            \mathcal{G}(U) \arrow[d, "\rho_{UV}"] \\
            \mathcal{F}(V) \arrow[r, "\varphi(V)"] &
            \mathcal{G}(V)
        \end{tikzcd}
    \end{equation*}
\end{definition}
\begin{remark}
    By property of direct limit, morphism $\varphi$ gives homomorphism $\varphi_{x}: \mathcal{F}_{x} \longrightarrow \mathcal{G}_{x}$ for all $x\in X$.
\end{remark}
\begin{definition}
    Let $X$ be a topological space, $\varphi: \mathcal{F} \longrightarrow \mathcal{G}$ morphism of (pre)sheaves of abelian groups on $X$.
    We say $\varphi$ is isomorphic if there exists morphism $\psi: \mathcal{G} \longrightarrow \mathcal{F}$ such that for each open subsets $U$, 
    $\psi(U)\varphi(U) = \identity_{\mathcal{F}(U)}$ and $\varphi(U)\psi(U) = \identity_{\mathcal{G}(U)}$.
\end{definition}
\begin{proposition}
    Let $X$ be a topological space, $\varphi: \mathcal{F} \longrightarrow \mathcal{G}$ morphism of sheaves of abelian groups on $X$.
    Then $\varphi$ is isomorphic if and only if $\varphi_{x}$ is isomorphic for all $x\in X$.
    \label{Proposition 1.2}
\end{proposition}
\begin{proof}
    "$\Rightarrow$": For all $(U, t)\in \mathcal{F}_{x}$, since $\varphi(V)$ is isomorphic,
    there exists $s\in \mathcal{F}(U)$ such that $\varphi(V)(s) = t$.
    Thus $\varphi_{x}((U, s)) = (U, t)$, get $\varphi_{x}$ is surjective.
    For all $(U, s)\in \ker(\varphi_{x})$, $(U, \varphi(U)(s)) = 0$ in $\mathcal{G}_{x}$.
    Then there exists open neighbourhood $V$ of $x$ such that $V \subseteq U$ and $\varphi(U)(s)\big{|}_{V} = 0$.
    By commutative diagram, $\varphi(U)(s)\big{|}_{V} = \varphi(V)(s\big{|}_{V})$.
    Since $\varphi(V)$ is isomorphic, get $s\big{|}_{V} = 0$ and so $(U, s) = 0$ in $\mathcal{F}_{x}$.
    Get $\varphi_{x}$ is isomorphic.
    \par
    "$\Leftarrow$": For each open subset $U$, want to show that $\varphi(U)$ is isomorphic.
    For $s\in \ker{\varphi(U)}$, $\varphi(U)(s) = 0$.
    Then for all $x\in U$, $(U, \varphi(U)(s)) = 0$ in $\mathcal{G}_{x}$.
    As $(U, \varphi(U)(s)) = \varphi_{x}((U, s))$ and $\varphi_{x}$ is isomorphic, get $(U, s) = 0$ in $\mathcal{F}_{x}$.
    Thus there exists open neighbourhood $W_{x}$ of $x$ such that $W_{x} \subseteq U$ and $s\big{|}_{W_{x}} = 0$.
    By property of sheaf, get $s = 0$ and so $\varphi(U)$ is injective.
    \par
    For all $t\in \mathcal{G}(U)$, $(U, t)$ is an element of $\mathcal{G}_{x}$ for all $x\in U$.
    Then there exists $(V_{x}, s(x))\in \mathcal{F}_{x}$ such that $\varphi_{x}((V_{x}, s(x))) = (U, t)$ and $\varphi(V_{x})(s(x)) = t\big{|}_{V_{x}}$.
    Note that $\varphi(V_{x} \cap V_{y})(s(x)\big{|}_{V_{x} \cap V_{y}}) = (t\big{|}_{V_{x}})\big{|}_{V_{x} \cap V_{y}} = (t\big{|}_{V_{y}})\big{|}_{V_{x} \cap V_{y}} = \varphi(V_{x} \cap V_{y})(s(y)\big{|}_{V_{x} \cap V_{y}})$.
    Since $\varphi(V_{x} \cap V_{y})$ is injective, get $s(x)\big{|}_{V_{x} \cap V_{y}} = s(y)\big{|}_{V_{x} \cap V_{y}}$.
    By property of sheaf, there exists $s\in \mathcal{F}(U)$ such that $s\big{|}_{V_{x}} = s(x)$.
    So $\varphi(U)(s)\big{|}_{V_{x}} = \varphi(V_{x})(s(x)) = t\big{|}_{V_{x}}$.
    By property of sheaf, get $\varphi(U)(s) = t$ and so $\varphi(U)$ is isomorphic.
\end{proof}
\begin{definition}
    Let $X$ be a topological space, $\varphi: \mathcal{F} \longrightarrow \mathcal{G}$ morphism of presheaves of abelian groups on $X$.
    We have several notions
    \begin{itemize}
        \item Define the kernel of $\varphi$ to be $\ker(\varphi): op(X) \longrightarrow Ab\quad U \longmapsto \ker(\varphi(U))$.
        \item Define the image of $\varphi$ to be $im(\varphi): op(X) \longrightarrow Ab\quad U \longmapsto im(\varphi(U))$.
        \item Define the cokernel of $\varphi$ to be $coker(\varphi): op(X) \longrightarrow Ab\quad U \longmapsto coker(\varphi(U))$.
    \end{itemize}
\end{definition}
\begin{remark}
    The definition of cokernel in fact induce the definition of quotient presheaf.
\end{remark}
\begin{definition}[\textbf{\emph{Sheafification}}]
    Let $X$ be a topological space, $\mathcal{F}$ presheaf of abelian groups on $X$.
    Define a sheaf $\mathcal{F}^{+}$ of abelian groups on $X$ by $U \longmapsto \mathcal{F}^{+}(U)$, 
    where $\mathcal{F}^{+}(U) = \{(s_{x})_{x\in U} \big{|} \text{ for all $x\in U$, $\exists$ \text{open} $V_{x}\ni x$ and $t(x)\in \mathcal{F}(V_{x})$ such that $t(x)_{y} = s_{y}$, $\forall y\in V_{x}$}\}$ is a subset of $\prod_{x\in U} \mathcal{F}_{x}$.
\end{definition}
\begin{proposition}
    Let $X$ be a topological space, $\mathcal{F}$ presheaf of abelian groups on X.
    Show that there is a natural morphism of presheaves $\varphi: \mathcal{F} \longrightarrow \mathcal{F}^{+}$ satisfies that
    \begin{itemize}
        \item $\varphi_{x}$ is identity homomorphism.
        \item If $\mathcal{F}$ is sheaf, then $\varphi$ is isomorphism.
        \item For any morphism of presheaves $\psi: \mathcal{F} \longrightarrow \mathcal{G}$ where $\mathcal{G}$ is a sheaf,
        there exists unique morphism $\phi: \mathcal{F}^{+} \longrightarrow \mathcal{G}$ such that the following diagram commutes.
        \begin{equation}
            \begin{tikzcd}
                \mathcal{F} \arrow[rr, "\psi"] \arrow[dr, "\varphi"] &&
                \mathcal{G} \\
                & \mathcal{F}^{+} \arrow[ur, "\phi"] &
            \end{tikzcd}
        \end{equation}
    \end{itemize}
    \label{Proposition 1.3}
\end{proposition}
\begin{reason}
    The construction of $\varphi$ is just $\varphi(U): s \longmapsto (s_{x})_{x\in U}$.
\end{reason}
\begin{definition}
    Let $X$ be a topological space, $\varphi: \mathcal{F} \longrightarrow \mathcal{G}$ morphism of sheaves of abelian groups on X. \\
    (1)Define the kernel of $\varphi$ to be $\ker(\varphi): op(X) \longrightarrow Ab\quad U \longmapsto \ker(\varphi(U))$. \\
    (2)Define the image of $\varphi$ to be sheafification of $op(X) \longrightarrow Ab\quad U \longmapsto im(\varphi(U))$, denoted by $im(\varphi)$. \\
    (3)Define the cokernel of $\varphi$ to be sheafification of $op(X) \longrightarrow Ab\quad U \longmapsto coker(\varphi(U))$, denoted by $coker(\varphi)$.
\end{definition}
\begin{remark}
    The definition of cokernel in fact induce the definition of quotient sheaf.
\end{remark}
\begin{definition}
    Let $X$ be a topological space, $\varphi: \mathcal{F} \longrightarrow \mathcal{G}$ morphism of sheaves of abelian groups on X. \\
    (1)We say that $\varphi$ is injective if $\ker(\varphi) = 0$. \\
    (2)We say that $\varphi$ is surjective if $im(\varphi) = \mathcal{G}$
\end{definition}
\begin{proposition}
    Let $X$ be a topological space, $\varphi: \mathcal{F} \longrightarrow \mathcal{G}$ morphism of sheaves of abelian groups on X.
    Then \\
    (1)$\varphi$ is injective if and only if $\varphi(U)$ is injective for all open subset U. \\
    (2)$\varphi$ is injective if and only if $\varphi_{x}$ is injective for all $x\in X$. \\
    (3)$\varphi$ is surjective if and only if $\varphi_{x}$ is surjective for all $x\in X$. \\
    (4)$\varphi$ is isomorphic if and only if $\varphi$ is injective and surjective.
    \label{Proposition 1.4}
\end{proposition}
\begin{remark}
    The proof is similar to proof of Proposition \ref{Proposition 1.2}.
    In addition, Hartshorne exercise II 1.2 gives a more explicit sketch.
\end{remark}
\begin{definition}[\textbf{\emph{Direct Image and Inverse Image}}]
    Let $f:X \longrightarrow Y$ be continuous map of topological spaces, $\mathcal{F}$ sheaf of abelian groups on X, $\mathcal{G}$ sheaf of abelian groups on Y. \\
    (1)Define dierct image to be the sheaf $f_{\ast}\mathcal{F}: V \longmapsto f_{\ast}\mathcal{F}(V) = \mathcal{F}(f^{-1}(V))$. \\
    (2)Define inverse image to be the sheafification of $U \longmapsto \underset{\underset{V \supseteq f(U)}{\longrightarrow}}{\lim} \mathcal{G}(V)$, denoted by $f^{-1}\mathcal{G}$.
\end{definition}
\begin{remark}
    In general, we don't have $(f_{\ast}\mathcal{F})_{f(x)} \neq \mathcal{F}_{x}$,
    while there is a canonical homomorphism $(f_{\ast}\mathcal{F})_{f(x)} \longrightarrow \mathcal{F}_{x}$.
    On the other hand, we always have $(f^{-1}\mathcal{G})_{x} = \mathcal{G}_{f(X)}$.
\end{remark}
\begin{example}
    Let $i: Z \hookrightarrow X$ be an inclusion of topological spaces, $\mathcal{F}$ sheaf of abelian groups on X.
    We often write $i^{-1}\mathcal{F}$ as $\mathcal{F}\big{|}_{Z}$.
\end{example}

\subsection{Sheaf modules}
\label{subsection:Vector_Bundles_Sheaf_modules}

\begin{definition}
    Let $X$ be a scheme or $(X, \mathcal{O}_{X})$ be a ringed space.
    A sheaf of $\mathcal{O}_{X}$-modules is a sheaf of abelian groups $\mathcal{F}$ on X such that for all $U \subseteq X$ open, 
    $\mathcal{F}(U)$ is a module over $\mathcal{O}_{X}(U)$ which is compatible with restriction maps.
\end{definition}
\begin{definition}
    Let $X$ be a scheme or $(X, \mathcal{O}_{X})$ be a ringed space, $\mathcal{F},\mathcal{G}$ two $\mathcal{O}_{X}$-modules.
    A morphism $\varphi$ from $\mathcal{F}$ to $\mathcal{G}$ is a morphism of sheaves compatible with the module structure i.e. for all $U \subseteq X$ open, 
    $\varphi(U): \mathcal{F}(U) \longrightarrow \mathcal{G}(U)$ is a homomorphism of $\mathcal{O}_{X}(U)$-modules.
\end{definition}
\begin{remark}
    Kernel, image and cokernel of a morphism of $\mathcal{O}_{X}$-modules are still $\mathcal{O}_{X}$-modules.
    Can also define submodule, quotient module and direct sum.
\end{remark}
\begin{definition}
    Let $X$ be a scheme or $(X, \mathcal{O}_{X})$ be a ringed space.
    An sheaf ideal on X is an $\mathcal{O}_{X}$-submodule of $\mathcal{O}_{X}$.
\end{definition}
\begin{definition}[\textbf{\emph{Tensor Products}}]
    Let $X$ be a scheme or $(X, \mathcal{O}_{X})$ be a ringed space, $\mathcal{F},\mathcal{G}$ two $\mathcal{O}_{X}$-modules.
    Define tensor product $\mathcal{F} \otimes_{\mathcal{O}_{X}} \mathcal{G}$ to be the sheafification of $U \longmapsto \mathcal{F}(U) \otimes_{\mathcal{O}_{X}(U)} \mathcal{G}(U)$.
\end{definition}
\begin{remark}
    The stalk of $\mathcal{F} \otimes_{\mathcal{O}_{X}} \mathcal{G}$ at $x\in X$ is just $\mathcal{F}_{x} \otimes_{\mathcal{O}_{X, x}} \mathcal{G}_{x}$.
\end{remark}
\begin{definition}[\textbf{\emph{Free $\mathcal{O}_{X}$-modules}}]
    Let $X$ be a scheme or $(X, \mathcal{O}_{X})$ be a ringed space, $\mathcal{F}$ an $\mathcal{O}_{X}$-module.
    We say that $\mathcal{F}$ is free if $\mathcal{F}$ is isomorphic to a direct sum of $\mathcal{O}_{X}$.
    In addition, we say that $\mathcal{F}$ is free of rank r if $\mathcal{F} \cong \mathcal{O}_{X}^{r}$.
\end{definition}
\begin{definition}[\textbf{\emph{Locally Free $\mathcal{O}_{X}$-modules}}]
    Let $X$ be a scheme or $(X, \mathcal{O}_{X})$ be a ringed space, $\mathcal{F}$ an $\mathcal{O}_{X}$-module.
    We say that $\mathcal{F}$ is locally free if X can be covered by open subsets $\{U_{i}\}_{i\in I}$ such that $\mathcal{F}\big{|}_{U_{i}}$ is free as $\mathcal{O}_{U_{i}}$-module.
    In addition, we say that $\mathcal{F}$ is locally free of rank r if X can be covered by open subsets $\{U_{i}\}_{i\in I}$ such that $\mathcal{F}\big{|}_{U_{i}} \cong \mathcal{O}_{U_{i}}^{r}$ as $\mathcal{O}_{U_{i}}$-module.
\end{definition}
\begin{definition}[\textbf{\emph{Invertible $\mathcal{O}_{X}$-modules}}]
    Let $X$ be a scheme or $(X, \mathcal{O}_{X})$ be a ringed space, $\mathcal{F}$ an $\mathcal{O}_{X}$-module.
    We say that $\mathcal{F}$ is invertible if it is locally free of rank 1.
\end{definition}
\begin{remark}
    If take $\mathcal{G} = \sheafhom_{\mathcal{O}_{X}}(\mathcal{F}, \mathcal{O}_{X})$, then $\mathcal{F} \otimes_{\mathcal{O}_{X}} \mathcal{G} \cong \mathcal{O}_{X}$.
\end{remark}
\begin{lemma}
    Let $f: X \longrightarrow Y$ be a morphism of schemes or ringed spaces, $\mathcal{F}$ an $\mathcal{O}_{X}$-module, $\mathcal{G}$ an $\mathcal{O}_{Y}$-module.
    Then the direct image $f_{\ast}\mathcal{F}$ of $\mathcal{F}$ is an $\mathcal{O}_{Y}$-module induced by $f^{\sharp}$.
    On the other hand, the inverse image $f^{\ast}\mathcal{G} = f^{-1}\mathcal{G} \otimes_{f^{-1}\mathcal{O}_{Y}} \mathcal{O}_{X}$ of $\mathcal{G}$ is an $\mathcal{O}_{X}$-module where $f^{-1}\mathcal{O}_{Y} \longrightarrow \mathcal{O}_{X}$ is induced by $f^{\sharp}$.
    \label{Lemma 1.1}
\end{lemma}

\subsection{Vector bundles}
\label{subsection:Vector_Bundles_Vector_bundles}

Let $X$ be an algebraic variety which we mean a separated scheme of finite type over $\mathbb{C}$.
And all morphisms mentioned are morphisms of varieties.
\begin{definition}
    A complex linear fibration over $X$ is a surjection $p: E \twoheadrightarrow X$ from variety $E$ such that for all $x\in X$, 
    the fiber $E_{x}$ is an affine space.
    In convention, we call $E$ the total space and $X$ the base variety.
\end{definition}
\begin{example}
    We call $X \times \mathbb{A}_{\mathbb{C}}^{r} \rightarrow X$ the trivial fibration of rank $r$.
\end{example}
\begin{lemma}
    Given a linear fibration $p: E \rightarrow X$, for any open subset $U \subseteq X$,
    the restriction $p^{-1}(U) \rightarrow U$ is also a linear fibration.
    Denote $E\big{|}_{U} = p^{-1}(U)$.
    \label{Lemma 1.2}
\end{lemma}
\begin{definition}
    Given two linear fibrations $E \overset{p}{\rightarrow} X$ and $F \overset{q}{\rightarrow} X$,
    a morphism $f: E \rightarrow F$ is said to be a morphism of linear fibrations if the following diagram commutes
    \begin{equation}
        \begin{tikzcd}
            E \arrow[rr, "f"] \arrow[dr, "p" swap] &&
            F \arrow[dl, "q"] \\
            & X &
        \end{tikzcd}
    \end{equation}
    and for all $x\in X$, $E_{x} \overset{f_{x}}{\rightarrow} F_{x}$ is a linear transformation.
\end{definition}
\begin{definition}
    An algebraic vector bundle of rank $r$ on $X$ is a linear fibration $E \overset{p}{\rightarrow} X$ which is locally trivial i.e. for all $x\in X$,
    there is an open neighbourhood $U$ of $x$ and isomorphism of linear fibrations $\varphi_{U}$ such that the following diagram commutes
    \begin{equation}
        \begin{tikzcd}
            E\big{|}_{U} \arrow[rr, "\varphi_{U}"] \arrow[dr, "p" swap] &&
            U \times \mathbb{A}_{\mathbb{C}}^{r} \arrow[dl] \\
            & U &
        \end{tikzcd}
    \end{equation}
    where $\varphi_{U}$ are called the local charts or local trivialization.
    \par
    And for two such $U_{i}$ and $U_{j}$ with local charts $\varphi_{i}$ and $\varphi_{j}$,
    there is an invertible matrix $g_{ij}$ with entries regular functions on $U_{ij}$, called transition function,
    satisfying that $g_{ii} = I_{r}$, $g_{ij} \cdot g_{jk} = g_{ik}$ and for all affine open subset $W \subseteq U_{ij}$,
    $\varphi_{i} \circ \varphi_{j}^{-1}\big{|}_{W}$ is linear transformation defined by $g_{ij}\big{|}_{W}$.
    Equivalently, $g_{ij}$ would correspond to a morphism $\widetilde{g}_{ij}: U_{ij} \rightarrow \generallineargroup{r}{\mathcal{O}_{X}(U_{ij})}$.
\end{definition}
\begin{lemma}
    Given data $\{g_{ij}\}$ satisfies that $g_{ii} = I_{r}$ and $g_{ij} \cdot g_{jk} = g_{ik}$, we can recover a vector bundle.
    \label{Lemma 1.3}
\end{lemma}
\begin{proof}
    Construct $E$ by gluing up $U_{i} \times \mathbb{A}_{\mathbb{C}}^{r}$ along $U_{ij} \times \mathbb{A}_{\mathbb{C}}^{r} \overset{\varphi_{ij}}{\rightarrow} U_{ij} \times \mathbb{A}_{\mathbb{C}}^{r}$ defined by $g_{ij}$.
\end{proof}
\begin{definition}
    Let $E, F$ be two vector bundles on $X$ of rank $r$ and $s$ respectively.
    We say a morphism $f: E \rightarrow F$ of linear fibrations is a morphism of vector bundles if there is a $r \times s$ matrix $g_{U}$ with entries regular functions on $U$ such that for all affine open subset $W \subseteq U$, 
    restriction of the following morphism to $W$ is linear transformation defined by $g_{U}\big{|}_{W}$
    \begin{equation}
        U \times \mathbb{A}_{\mathbb{C}}^{r} \overset{\varphi_{U}^{-1}}{\rightarrow} E\big{|}_{U} \overset{f\big{|}_{U}}{\rightarrow} F\big{|}_{U} \overset{\psi_{U}}{\rightarrow} U \times \mathbb{A}_{\mathbb{C}}^{s}
    \end{equation}
    where $\varphi_{U}$ and $\psi_{U}$ are local charts.
    Equivalently, $g_{U}$ would correspond to a morphism $\widetilde{g}_{U}: U \rightarrow \generallineargroup{r}{\mathcal{O}_{X}(U)}$.
\end{definition}
\begin{remark}
    With Lemma \ref{Lemma 1.3}, we can easily define $E \oplus F$, $E \otimes F$, $\sheafhom(E, F)$, $E^{\vee}$, $\wedge^{k} E$ and $S^{k} E$.
\end{remark}
\begin{proposition}
    Let $f: E \rightarrow F$ be a morphism of vector bundles.
    Assume that for $x_{0}\in X$, $f_{x_{0}}: E_{x_{0}} \rightarrow F_{x_{0}}$ is invertible.
    Then there exists an open neighbourhood $U_{0}$ of $x_{0}$ such that $f\big{|}_{U}: E\big{|}_{U} \rightarrow F\big{|}_{U}$ is an isomorphism.
    \label{Proposition 1.5}
\end{proposition}
\begin{proof}
    Assume $x_{0}\in U$ with local charts $\varphi_{U}: E\big{|}_{U} \rightarrow U \times \mathbb{A}_{\mathbb{C}}^{r}$ and $\psi_{U}: F\big{|}_{U} \rightarrow \mathbb{A}_{\mathbb{C}}^{s}$.
    By definition, there is a matrix $g_{U}$ and $f_{x_{0}}$ is defined by $g_{U}(x_{0})$.
    This immediately comes from the fact that there exists an open neighbourhood $U_{0} \subseteq U$ of $x_{0}$ such that $g_{U}\big{|}_{U_{0}}$ is invertible.
\end{proof}
\begin{definition}
    Let $E$ be a vector bundle of rank $r$ on $X$.
    A subbundle of rank $m$ of $E$ is a subvariety $F \subseteq E$ such that fiber $F_{x}$ is an $m$-dimensional affine linear subspace of $E_{x}$ and the induced fibration $p\big{|}_{F}: F \rightarrow X$ is locally trivial.
\end{definition}
\begin{example}
    Let $X = \mathbb{P}_{\mathbb{C}}^{n}$ be $n$-dimensional projective space with homogeneous coordinate $x_{0}, \cdots, x_{n}$.
    Consider the subvariety $\Sigma := V(x_{i}y_{j} - x_{j}y_{i})$ of $X \times \mathbb{A}_{\mathbb{C}}^{n + 1}$,
    where $y_{0}, \cdots, y_{n}$ are coordinate of $\mathbb{A}_{\mathbb{C}}^{n + 1}$.
    Then the fiber is a $1$-dimensional affine space and hence $\Sigma$ is a subbundle of rank $1$ of the trivial bundle, denoted by $\mathcal{O}(-1)$.
\end{example}
\begin{proposition}
    Let $F \subseteq E$ be a subbundle of rank $m$.
    Then there exists unique vector bundle structure on $E/F$ satisfying the universal property that any map of vector bundles $E \rightarrow G$ vanishing over $F$ would factor through $E/F$ uniquely.
    \label{Proposition 1.6}
\end{proposition}
\begin{proof}
    Firstly cover $X$ with affine open subsets trivializing $F$ and $E$, we may assume that $X = \spec A$, $E = \mathbb{A}_{A}^{n}$ and $F = \mathbb{A}_{A}^{m}$.
    Since $F \subseteq E$ is a subbundle, by definition, for all $x\in X$,
    $E_{x} = \mathbb{A}_{k(x)}^{n}$ and $F_{x} = \mathbb{A}_{k(x)}^{m}$ with coordinates $x_{1}, \cdots, x_{n}$ and $y_{1}, \cdots, y_{m}$ respectively,
    satisfying that $y_{i}$ can be linearly represented by $x_{i}$.
    \par
    It is clear that these relations still hold over a neighbourhood $U = D(f)$ of $x$.
    Then we extend $y_{i}$ to a basis $y_{1}, \cdots, y_{m}, z_{1}, \cdots, z_{n - m}$ and hence $\spec A_{f}[x_{1}, \cdots, x_{n}]/(y_{1}, \cdots, y_{m}) \cong \spec A_{f}[z_{1}, \cdots, z_{n - m}]$ is a linear subspace.
    For different $D(f)$ and $D(g)$, since $z_{1}, \cdots, z_{n - m}$ and $z_{1}', \cdots, z_{n - m}'$ span same linear subspace,
    there is an invertible transformation matrix $M\in \matrices{n - m}{A_{fg}}$ i.e. the transition function.
    Thus we defines a vector bundle structure on $E/F$.
    \par
    Now given a map of vector bundles $E \rightarrow G$ vanishing over $F$.
    Still assume that $X = \spec A$ trivializing all those bundles.
    Then $E \rightarrow G$ corresponds to a linear transformation $\varphi: A[w_{1}, \cdots, w_{r}] \rightarrow A[x_{1}, \cdots, x_{n}]$,
    which is defined by an $n \times r$ matrix with entries in $A$.
    By assumption, we know $y_{i}$ are not in the image of $\varphi$ and hence $\im \varphi$ lies in the subspace spanned by $z_{i}$.
    Thus $\varphi$ factors a canonical map $A[w_{1}, \cdots, w_{r}] \rightarrow A[z_{1}, \cdots, z_{n - m}]$.
    Finally, by universal property, we know the vector bundle structure on $E/F$ is also unique.
\end{proof}
\begin{proposition}[\textbf{\emph{Rigidity Property}}]
    Let $f: E \rightarrow F$ be a map of vector bundles.
    Suppose that rank of fiber morphisms $f_{x}$ is constant as $x$ varies over $X$.
    Then $\ker f$ and $\im f$ are subbundles of $E$ and $F$ respectively.
    \label{Proposition 1.7}
\end{proposition}
\begin{proof}
    Here we only show $\ker f$ is a subbundle of $E$ and similar argument makes sense for $\im f$.
    We may assume that $X = \spec A$ trivializing both $E$ and $F$.
    Now $f$ corresponds to a linear transformation $\varphi: A[y_{1}, \cdots, y_{m}] \rightarrow A[x_{1}, \cdots, x_{n}]$ defined by an $n \times m$ matrix $M$ with entries in $A$.
    \par
    Since the rank of fiber morphism is constant, we know both the dimensions of $\im \varphi_{x}$ are constant, denoted by $n - r$.
    Then we may choose $y_{1}, \cdots, y_{r}$ spanning the complementary subspace of $\im \varphi_{x}$, which can be linearly represented by $x_{i}$.
    And these relation still hold over an neighbourhood $U = D(f)$ of $x$.
    For different $D(f)$ and $D(g)$, since both $y_{1}, \cdots, y_{r}$ and $y_{1}', \cdots, y_{r}'$ span complementary subspace of $\im \varphi_{fg}$,
    there is an invertible transformation matrix $M\in \matrices{r}{A_{fg}}$ i.e. the transition function.
    Thus we give a subbundle structure on $\ker f$.
\end{proof}
\begin{remark}
    Though proof for $\im f$ follows by similar argument, there is still a little tricky thing in that case.
    When we talk about $\ker f_{x}$, it is clear the points living outside $\im \varphi_{x}$.
    But when we talk about $\im f_{x}$, what are they?
    As $X$ is locally noetherian, by a proposition in Some Useful Results, we know they are just the points living in coimage of $\varphi_{x}$.
    \par
    In addition, since this note is translated by me from language of varieties to language of schemes,
    there must be some ambiguity about the terminology.
    For example, here by $\ker f$, we possibly mean the subset of points whoes image under $f$ are same to image of the origin.
\end{remark}
\begin{corollary}
    Let $f: E \rightarrow F$ be a map of vector bundles.
    When $f$ is surjective, $\ker f$ is a subbundle of $E$.
    And when $f$ is injective, $\im f$ is a subbundle of $F$.
    \label{Corollary 1.1}
\end{corollary}
\begin{reason}
    When $f$ is surjective, as remark above, we know $\ker \varphi = 0$ so that rank of $f_{x}$ is constant.
    When $f$ is injective, as proof above, we know $\im \varphi = k[x_{1}, \cdots, x_{n}]$ so that rank of $f_{x}$ is constant.
\end{reason}
\begin{definition}
    Let $p: E \rightarrow X$ be a vector bundle on $X$ of rank $r$.
    A (regular) section of $E$ over open subset $U \subseteq X$ is a morphism $s: U \rightarrow E$ such that $(p \circ s)\big{|}_{U} = \identity_{U}$.
    Denote $\Gamma(U, E) = \{\text{regular sections of $E$ over $U$}\}$.
\end{definition}
\begin{lemma}
    Let $p: E \rightarrow X$ be a vector bundle on $X$ of rank $r$.
    Then $\Gamma(U, E)$ naturally admits an $\mathcal{O}(U)$-module structure
    \begin{equation}
        \mathcal{O}(U) \times \Gamma(U, E) \longrightarrow \Gamma(U, E)
    \end{equation}
    Moreover, $\mathcal{O}_{X}(E)(U) = \Gamma(U, E)$ gives a locally free sheaf $\mathcal{O}_{X}(E)$ of rank $r$.
    \label{Lemma 1.4}
\end{lemma}
\begin{proof}
    For $\alpha\in \mathcal{O}(U)$ and section $s: U \rightarrow E$, locally $s$ is defined by a homomorphism $\varphi: A[x_{1}, \cdots, x_{r}] \rightarrow A$.
    Define $\alpha \cdot s$ by gluing up $\alpha \cdot \varphi: x_{i} \mapsto \alpha \cdot \varphi(x_{i})$.
    In particular, locally $\Gamma(U, E) \cong A^{\oplus r}$ hence $\mathcal{O}_{X}(E)$ is a locally free sheaf of rank $r$. 
\end{proof}
\begin{proposition}
    Let $\mathcal{C}$ be the category of locally free sheaves of finite rank.
    There is a functor $\theta: \vectorbundle_{X} \rightarrow \mathcal{C}$ sending $E$ to $\mathcal{O}_{X}(E)$ and $E \rightarrow F$ to $\{\Gamma(U, E) \rightarrow \Gamma(U, F)\}$,
    which is an equivalence of category.
    \label{Proposition 1.8}
\end{proposition}
\begin{proof}
    First to prove $\theta$ fully faithful i.e. $\hom_{\vectorbundle_{X}}(E, F) \rightarrow \hom_{\mathcal{O}_{X}}(\mathcal{O}_{X}(E), \mathcal{O}_{X}(F))$ is bijection for all $E, F$.
    We may assume that $X = \spec A$ and $E, F$ are both trivial.
    Then $f: \mathbb{A}_{A}^{n} \rightarrow \mathbb{A}_{A}^{m}$ corresponds to a matrix $M\in \matrices{n \times m}{A}$.
    Also, $(s: U \rightarrow E) \rightarrow (f \circ s: U \rightarrow F)$ also corresponds to such a matrix.
    Hence the bijection is clear.
    \par
    Next to prover $\theta$ dense (or say essentially surjective) i.e. given $\mathcal{E}\in \mathcal{C}$, there exists $E\in \vectorbundle_{X}$ with isomorphism $\mathcal{O}_{X}(E) \overset{\sim}{\rightarrow} \mathcal{E}$.
    Assume $U_{i} = \spec A_{i}$ is an affine cover of $X$ trivializing $\mathcal{E}$.
    Denote the trivialization $\varphi_{i}: \mathcal{E}\big{|}_{U_{i}} \rightarrow \mathbb{A}_{A_{i}}^{n}$.
    Consider the transition functions $g_{ij} = \varphi_{i} \circ \varphi_{j}^{-1}\big{|}_{U_{ij}}$.
    Then thoes transition functions give a vector bundle $E$.
    \par
    Define $\psi_{i}: \mathcal{O}_{X}(E)(U_{i}) = \Gamma(U_{i}, E) \rightarrow \mathcal{E}\big{|}_{U_{i}}$ sending $s: U_{i} \rightarrow \mathbb{A}_{U_{i}}^{n}$ to $\varphi_{i}^{-1}(v_{s})$,
    where $v_{s}$ is the vector defining $s$.
    Then 
    \begin{equation}
        \psi_{i}\big{|}_{U_{ij}}(s\big{|}_{U_{ij}}) = \varphi_{i}^{-1}(v_{s}) = \varphi_{i}^{-1}(g_{ij}v_{s}') = \varphi_{i}^{-1} \circ \varphi_{i} \circ \varphi_{j}^{-1}(v_{s}') = \varphi_{j}^{-1}(v_{s}') = \psi_{j}\big{|}_{U_{ij}}(s\big{|}_{U_{ij}})
    \end{equation}
    Hence $\psi_{i}$ glue up and get our desired isomorphism.
\end{proof}

\subsection{Bundles associated to a representation}
\label{subsection:Vector_Bundles_Bundles_associated_to_a_representation}

\begin{lemma}
    Let $f: Y \rightarrow X$ be a morphism, $p: E \rightarrow X$ vector bundle of rank $r$ on $X$.
    Consider the fibered product $F = Y \times_{X} E$.
    Then the first projection $q: F \rightarrow Y$ gives a bundle on $Y$ of rank $r$ and $\mathcal{O}_{Y}(F) = f^{\ast}\mathcal{O}_{X}(E)$.
    \label{Lemma 1.5}
\end{lemma}
\begin{proof}
    Locally we have trivializations $\varphi_{i}: E\big{|}_{U_{i}} \rightarrow U_{i} \times \mathbb{A}_{\mathbb{{C}}}^{r}$ and transition functions $g_{ij}: U_{ij} \times \mathbb{A}_{\mathbb{{C}}}^{r} \rightarrow U_{ij} \times \mathbb{A}_{\mathbb{{C}}}^{r}$ satisfying cocycle condition.
    Then after base change we get $\widetilde{\varphi}_{i}: F\big{|}_{V_{i}} \rightarrow V_{i} \times \mathbb{A}_{\mathbb{{C}}}^{r}$ and $\widetilde{g}_{ij}: U_{ij} \times \mathbb{A}_{\mathbb{{C}}}^{r} \rightarrow U_{ij} \times \mathbb{A}_{\mathbb{{C}}}^{r}$ still satisfying cocycle condition.
    Hence $q: F \rightarrow Y$ defines a bundle on $Y$ of rank $r$.
    \par
    Now for all open subset $V \subseteq Y$, define $\sigma(V): f^{\ast}\mathcal{O}_{X}(E)(V) \rightarrow \mathcal{O}_{Y}(F)(V)$ sending $(U, s) \otimes 1$ to $\widetilde{s}\big{|}_{V}$,
    where $\widetilde{s}: f^{-1}(U) \rightarrow F$ is the pull back of $s: U \rightarrow E$.
    For all $y\in Y$, want to show that $\sigma_{y}: \mathcal{O}_{X}(E)_{f(y)} \otimes_{\mathcal{O}_{X, f(y)}} \mathcal{O}_{Y, y} \rightarrow \mathcal{O}_{Y}(F)_{y}$ is isomorphism.
    \par
    For injectivity, if $\sum_{i = 1}^{r} (U, s_{i}) \otimes b_{i}\in \ker \sigma_{y}$,
    then locally we get $\sum_{i = 1}^{r} b_{i} \cdot f^{\sharp} \circ s_{i}^{\sharp}(x_{j}) = 0$ for all $1 \le j \le r$,
    so that also $\sum_{i = 1}^{r} (s_{i}^{\sharp}(x_{1}), \cdots, s_{i}^{\sharp}(x_{r})) \otimes b_{i} = 1 \otimes (\sum_{i = 1}^{r} b_{i} \cdot f^{\sharp} \circ s_{i}^{\sharp}(x_{1}), \cdots, \sum_{i = 1}^{r} b_{i} \cdot f^{\sharp} \circ s_{i}^{\sharp}(x_{r})) = 0$.
    Then $(U, s) = \oplus_{i = 1}^{r} a_{i} \otimes 1 = 0$.
    For surjectivity, given $(b_{1}, \cdots, b_{r})\in B^{\oplus r}$,
    clearly it is the image of $\sum_{i = 1}^{r} (U, s_{i}) \otimes b_{i}$ where $s_{i}$ is defined by $e_{i}\in A^{\oplus r}$.
\end{proof}
Let $p: E \rightarrow X$ be a vector bundle of rank $r$.
To avoid ambiguity, we denote the general linear group by $G_{r}$ and the group variety by $\generallineargroup{r}{\mathbb{C}}$,
which is the open subset $D(\det)$ of $\mathbb{A}_{\mathbb{C}}^{r^{2}}$.
Consider the variety $R = X \times \generallineargroup{r}{\mathbb{C}}$.
We have a group action of $G_{r}$ on $R$ defined locally as following, for $M = (a_{ij})$
\begin{equation}
    \begin{split}
        \identity \times (\circ M^{-1}): A \times \generallineargroup{r}{\mathbb{C}} \rightarrow A \times \generallineargroup{r}{\mathbb{C}}
    \end{split}
\end{equation}
where $\circ M^{-1}$ corresponds to the homomorphism defined by $(x_{ij}) \mapsto (\sum_{k} x_{ik}a_{kj})$.
Then $R/G_{r} = X$ and the orbits can be identified with the fibers of $R \rightarrow X$.
In addition, each matirx $M$ gives a morphism $\mathbb{A}_{\mathbb{C}}^{r} \rightarrow E\big{|}_{U}$ satisfying that composition with trivialization $\varphi_{U}$ is a linear transformation defined by $M$.
Then it is natural to consider transition function $g_{ij} \otimes I_{r}$ which gives a $\generallineargroup{r}{\mathbb{C}}$-bundle on $X$, denoted by $\framebundle (E)$, the frame bundle of $E$.
And there is an group action of $G_{r}$ on $R \times \mathbb{A}_{\mathbb{C}}^{r}$ defined locally as following, for $M = (a_{ij})$
\begin{equation}
    \begin{split}
        \identity \times (\circ M^{-1}) \times M: A \times \generallineargroup{r}{\mathbb{C}} \times \mathbb{A}_{\mathbb{C}}^{r} \rightarrow A \times \generallineargroup{r}{\mathbb{C}} \times \mathbb{A}_{\mathbb{C}}^{r}
    \end{split}
\end{equation}
where $M$ corresponds to the linear transformation defined by $M$.
\par
Now given a representation $\psi: G_{r} \rightarrow G_{m}$,
consider the group action of $G_{r}$ on $R \times \mathbb{A}_{\mathbb{C}}^{m}$ defined locally as following, for $M = (a_{ij})$
\begin{equation}
    \begin{split}
        \identity \times (\circ M^{-1}) \times \psi(M): A \times \generallineargroup{r}{\mathbb{C}} \times \mathbb{A}_{\mathbb{C}}^{m} \rightarrow A \times \generallineargroup{r}{\mathbb{C}} \times \mathbb{A}_{\mathbb{C}}^{m}
    \end{split}
\end{equation}
Then locally the quotient scheme is isomorphic to $\mathbb{A}_{A}^{m}$ with natural transition functions.
These information give a vector bundle $E_{\psi}$ of rank $m$ on $X$, called the bundle associated to representation $\psi$.
\begin{remark}
    The idea of frame bundle also comes from moduli thery.
    Consider functor $\mathcal{F}: \scheme_{X} \rightarrow \sets$ defined as following
    \begin{equation}
        (f: T \rightarrow X) \longmapsto \{f^{\ast}\mathcal{O}(E) \overset{\sim}{\rightarrow} \mathcal{O}_{T}^{r}\}
    \end{equation}
    And in fact this functor is representable by frame bundle above.
\end{remark}

\section{Riemann-Roch Formula}
\label{section:Riemann-Roch_Formula}
\end{document}